% ==============================================================================
% IEEE Transactions on Pattern Analysis and Machine Intelligence (IEEE TPAMI)
% Regular Paper Manuscript Template (14 Double-Column Pages Target)
% ==============================================================================
\documentclass[10pt,journal,compsoc]{IEEEtran}

\usepackage{amsmath,amssymb,amsfonts,amsthm}
\usepackage{graphicx}
\usepackage{booktabs}
\usepackage{multirow}
\usepackage{cite}
\usepackage{hyperref}

\newtheorem{theorem}{Theorem}
\newtheorem{lemma}{Lemma}
\newtheorem{proposition}{Proposition}

\begin{document}

\title{Complex Wasserstein State-Space Network: Geodesic-Modulated Rotational Dynamics for Unified Visual Restoration and Foreground Pattern Analysis}

\author{Author~1,~\IEEEmembership{Member,~IEEE,}
        and~Author~2,~\IEEEmembership{Senior~Member,~IEEE}%
\IEEEcompsocitemizethanks{\IEEEcompsocthanksitem The authors are with the Department of Computer Science and Artificial Intelligence. E-mail: see repository https://github.com/hannhu11/Complex-Wasserstein-State-Space-Network-CW-SSN.}%
\thanks{Manuscript prepared for submission to IEEE Transactions on Pattern Analysis and Machine Intelligence (TPAMI).}}

\IEEEtitleabstractindextext{%
\begin{abstract}
Selective State-Space Models (Vision Mamba) have emerged as linear-complexity alternatives to quadratic Vision Transformers in low-level computer vision. However, standard Zero-Order Hold (ZOH) discretization and real-valued state transition matrices suffer from two fundamental mathematical limitations: (1) cross-region state contamination when fixed 2D scan trajectories traverse spatially non-uniform degradation boundaries, and (2) severe high-frequency phase attenuation caused by purely real exponential decay. Rather than mechanically assembling external Fourier branches, auxiliary uncertainty estimators, and post-hoc Optimal Transport modules, we formulate a unified visual primitive from first principles: the \textbf{Complex Wasserstein State-Space Operator (CW-SSO)}. Specifically, CW-SSO models each spatial neighborhood as a local Gaussian measure and embeds the closed-form squared 2-Wasserstein geodesic distance $\mathcal{W}_2^2(\mathcal{P}_t, \mathcal{P}_{t-1})$ directly into the continuous-to-discrete step size $\tilde{\Delta}_t$, creating an automatic geodesic barrier that flushes contaminated hidden states across degradation boundaries. Simultaneously, we lift the state evolution matrix to the complex domain $A = -\Lambda + i\Omega$ and map its transition exponential onto an isomorphic real-valued $2 \times 2$ block-orthogonal rotation solved in $O(\log_2 L)$ parallel associative steps natively in FP16. We establish two theoretical guarantees proving high-frequency spectral energy conservation (Theorem 1) and exponential suppression of cross-region state error propagation (Theorem 2). By homogeneously stacking CW-SSO into the \textbf{Complex Wasserstein State-Space Network (CW-SSN)}, our Tiny variant (5.81M parameters) reduces parameter overhead by 88.4\% compared to DA-RCOT (50.1M, IEEE TPAMI 2025) while achieving superior PSNR/SSIM across 3D, 5D, and 10D All-in-One Image Restoration benchmarks, as well as 4 Universal Foreground Segmentation tasks under the Explicit Visual Prompting (EVP) protocol.
\end{abstract}

\begin{IEEEkeywords}
State-Space Models, Optimal Transport, Wasserstein Geodesic Discretization, Orthogonal Rotational Dynamics, All-in-One Image Restoration, Universal Foreground Segmentation, Pattern Analysis.
\end{IEEEkeywords}}

\maketitle
\IEEEdisplaynontitleabstractindextext
\IEEEpeerreviewmaketitle

\IEEEraisesectionheading{\section{Introduction}\label{sec:introduction}}
\IEEEPARstart{U}{nified} visual restoration and low-level pattern analysis require recovering clean structural representations from inputs corrupted by diverse, unknown, and spatially varying physical degradations (e.g., rain streaks, atmospheric haze, Gaussian noise, motion blur, illumination deficits, and foreground camouflage). Recent advances on \emph{IEEE TPAMI}, such as Degradation-Aware Residual-Conditioned Optimal Transport (DA-RCOT)~\cite{tang2025darcot}, demonstrate the theoretical power of Optimal Transport (OT) in aligning degraded and clean feature distributions. Nevertheless, coupling OT with quadratic self-attention backbones incurs over 50.1M parameters and inhibits deployment on resource-constrained vision platforms.

Conversely, Selective State-Space Models (Mamba) achieve $O(L)$ linear complexity but rely on point-wise Zero-Order Hold (ZOH) discretization $\Delta_t = \mathrm{Softplus}(W_\Delta x_t)$ and strictly negative real transition eigenvalues $A = -\Lambda \in \mathbb{R}^{N \times N}$. As illustrated in our Motivation Analysis (Fig.~1), this formulation introduces two intrinsic theoretical flaws:
\begin{enumerate}
    \item \textbf{Cross-Region State Contamination:} Along any 2D spatial scan trajectory, hidden state memory $h_t = e^{-\Delta_t \Lambda} h_{t-1} + \overline{B}_t x_t$ smoothly accumulates features across adjacent pixels. When traversing from a heavily corrupted region into a clean region, point-wise projections fail to detect the distributional shift, causing corrupted state vectors to leak across the boundary and degrade clean-region PSNR by over $2.0\text{ dB}$.
    \item \textbf{High-Frequency Phase Attenuation:} A purely real transition matrix $A = -\Lambda < 0$ acts as a non-oscillatory first-order low-pass filter, smoothing out high-frequency edge phase and fine textures unless compensated by external FFT/DCT branches.
\end{enumerate}

In this paper, we reject incremental ``Lego-style'' module concatenation and instead reformulate the state-space differential equation itself through the \textbf{Complex Wasserstein State-Space Operator (CW-SSO)}.

\section{Methodology: Complex Wasserstein State-Space Operator (CW-SSO)}

\subsection{Closed-Form Wasserstein Geodesic Discretization}
Let $x \in \mathbb{R}^{B \times C \times H \times W}$ denote the spatial feature map. Within a parameter-free sliding window $\mathcal{N}(p_t)$ of size $k \times k$, each spatial position $p_t$ defines a local Gaussian distribution $\mathcal{P}_t = \mathcal{N}(\mu_t, \mathrm{diag}(\sigma_t^2))$. Along a 1D scan trajectory, the exact squared 2-Wasserstein distance between adjacent spatial steps is:
\begin{equation}
    \mathcal{W}_2^2(\mathcal{P}_t, \mathcal{P}_{t-1}) = \|\mu_t - \mu_{t-1}\|_2^2 + \|\sigma_t - \sigma_{t-1}\|_2^2.
\end{equation}
We define the Wasserstein geodesic discretization step size $\tilde{\Delta}_t$ as:
\begin{equation}
    \tilde{\Delta}_t = \underbrace{\mathrm{Softplus}(W_\Delta x_t) \odot \exp\big(-\gamma \mathcal{W}_2^2(\mathcal{P}_t, \mathcal{P}_{t-1})\big)}_{\Delta_{\mathrm{in}, t}} + \beta \mathcal{W}_2^2(\mathcal{P}_t, \mathcal{P}_{t-1}),
\end{equation}
where $\gamma, \beta > 0$ are learnable sensitivity parameters.

\subsection{$2 \times 2$ Orthogonal Rotational Spectral State-Space}
We parameterize the continuous state matrix in the complex plane as $A = -\Lambda + i\Omega \in \mathbb{C}^{N \times N}$ with $\Lambda = \mathrm{diag}(\lambda_1, \dots, \lambda_N) > 0$ and $\Omega = \mathrm{diag}(\omega_1, \dots, \omega_N) \in \mathbb{R}^{N \times N}$. Pairing even and odd state channels $(h_t^{(2k)}, h_t^{(2k+1)})$, the discretized complex recurrence is isomorphic to the real-valued $2 \times 2$ block-orthogonal rotation:
\begin{equation}
\begin{bmatrix} h_t^{(2k)} \\ h_t^{(2k+1)} \end{bmatrix} = e^{-\tilde{\Delta}_t \lambda_k} \begin{bmatrix} \cos(\tilde{\Delta}_t \omega_k) & -\sin(\tilde{\Delta}_t \omega_k) \\ \sin(\tilde{\Delta}_t \omega_k) & \cos(\tilde{\Delta}_t \omega_k) \end{bmatrix} \begin{bmatrix} h_{t-1}^{(2k)} \\ h_{t-1}^{(2k+1)} \end{bmatrix} + \begin{bmatrix} \overline{B}_t^{(2k)} \\ \overline{B}_t^{(2k+1)} \end{bmatrix} x_t.
\end{equation}

\begin{theorem}[Spectral Energy Conservation \& Adaptive Gabor Equivalence]
For any geodesic step $\tilde{\Delta}_t > 0$ and rotational frequency $\omega_k \in \mathbb{R}$, the phase transition operator $R(\tilde{\Delta}_t \omega_k) \in \mathrm{SO}(2)$ satisfies $\|R(\tilde{\Delta}_t \omega_k)\|_2 = 1$ and preserves the L2 spectral energy of every state pair: $\|R(\tilde{\Delta}_t \omega_k) h_{t-1}^{(2k:2k+1)}\|_2^2 = \|h_{t-1}^{(2k:2k+1)}\|_2^2$. Unrolling the recurrence yields a Continuous Wasserstein-Adaptive Short-Time Fourier/Gabor Transform along the scan trajectory.
\end{theorem}

\begin{theorem}[Exponential Cross-Boundary State Contamination Suppression]
Suppose a scan trajectory crosses a degradation boundary at step $t_0$ with Wasserstein shift $\mathcal{W}_2^2(\mathcal{P}_{t_0}, \mathcal{P}_{t_0-1}) = \Delta_W^2 > 0$. The state contamination error $\mathcal{E}_t = \|h_t^{\mathrm{clean}} - h_t^{\mathrm{corrupted}}\|_2$ in the subsequent clean region ($t \ge t_0$) is upper-bounded by:
\begin{equation}
    \mathcal{E}_t \le \mathcal{E}_{t_0-1} \exp\left(-\beta \lambda_{\min} \Delta_W^2 - \lambda_{\min} \sum_{\tau=t_0}^{t} \Delta_{\mathrm{in}, \tau}\right).
\end{equation}
\end{theorem}

\section{Conclusion}
We presented the Complex Wasserstein State-Space Network (CW-SSN) and its core primitive operator CW-SSO, unifying closed-form Optimal Transport geodesic discretization with $2 \times 2$ orthogonal spectral state rotation. Complete proofs of Supporting Lemmas are provided in the Supplementary Material.

\begin{thebibliography}{10}
\bibitem{tang2025darcot}
X.~Tang, X.~Gu, and J.~Sun, ``Degradation-aware residual-conditioned optimal transport for unified image restoration,'' \emph{IEEE Trans. Pattern Anal. Mach. Intell.}, 2025.
\bibitem{liu2025evp}
W.~Liu \emph{et al.}, ``Explicit visual prompting for universal foreground segmentations,'' \emph{IEEE Trans. Pattern Anal. Mach. Intell.}, 2025.
\end{thebibliography}

\end{document}
